% user_guide.tex — complete reference for every user-facing feature
\chapter{User Guide}
\label{chap:user-guide}

This chapter is the authoritative reference for every user-facing command,
configuration option, document format, and error condition exposed by
\cgspkg{}.

% -----------------------------------------------------------------------
\section{CLI Overview}

The single entry-point is the \cgscli{} command.

\begin{lstlisting}[language=bash]
$ cgitsync <command> [options]
\end{lstlisting}

Running \cgscli{} without a command prints a help summary and exits with
code~0.

% -----------------------------------------------------------------------
\section{Lifecycle Contract}

This guide follows the simplified canonical lifecycle:

\begin{enumerate}
  \item \texttt{initialise(.cgs)} $\rightarrow$ clone all repos $\rightarrow$
        \texttt{READY}  \textit{(new project)}\\
        \texttt{initialise(.gts)} $\rightarrow$ restore from snapshot $\rightarrow$
        \texttt{READY}  \textit{(existing project)}
  \item \texttt{pull(.cgs|.gts)} $\rightarrow$ resync an existing tree
  \item \texttt{checkout(.gts)} / \texttt{add} / \texttt{commit} / \texttt{push} / \texttt{tag}
        $\rightarrow$ git operations on the tree
  \item \texttt{freeze} $\rightarrow$ emit the next \texttt{.gts} id and update \texttt{<Project\_name>.lgr}
\end{enumerate}

Compatibility commands remain available for older scripts:
\texttt{clone} can still initialise from a \texttt{.cgs}, \texttt{restart}
is an alias-style resynchronisation command, and \texttt{freeze-release} /
\texttt{freeze-state} keep their specialised names. New CLI usage should start
with \texttt{initialise}.

The project-local \texttt{.lgr} register and sync ledger are both active (T24,
T32).  Every synchronisation operation is automatically recorded as an immutable
DAG event in the \texttt{[[ledger]]} section of the \texttt{.lgr} file.

% -----------------------------------------------------------------------
\section{Commands}

\subsection{\texttt{initialise}}

\begin{lstlisting}[language=bash]
$ cgitsync initialise <spec.cgs> [--output-path DIR]
$ cgitsync initialise <snapshot.gts>
\end{lstlisting}

Unified initialisation entry point.  Dispatches on file extension:

\begin{itemize}
  \item \texttt{.cgs} source (clone mode): clones the full repository tree
    (see \texttt{clone} below). The optional \texttt{--output-path} controls
    \texttt{CGSPATH}; \texttt{CGSHOME} is derived as
    \texttt{CGSPATH/<project-name>}. When omitted, \texttt{CGSPATH} defaults to
    \texttt{../..}.  On success, ends in \texttt{READY}.  If clone setup fails
    because stale generated state is still present, the CLI prints
    \texttt{Try clean-init method}.
  \item \texttt{.gts} source (restore mode): loads a saved snapshot directly.
    On success, ends in the lifecycle state recorded in the snapshot.
\end{itemize}

Both paths end in a \texttt{READY} tree or fail explicitly.

\subsection{\texttt{clean-init}}

\begin{lstlisting}[language=bash]
$ cgitsync clean-init <spec.cgs> [--output-path DIR]
\end{lstlisting}

Recovery-oriented initialisation for a \texttt{.cgs} source. It follows the
same high-level pipeline as \texttt{initialise(.cgs)}, but inserts a cleanup
phase before cloning:

\begin{center}
\texttt{load->expand->validate->purge->clone}
\end{center}

Use it when an earlier initialisation attempt left partial child checkouts,
stale submodule metadata, or an obsolete root \texttt{.gitmodules} file.

\subsection{\texttt{purge}}

\begin{lstlisting}[language=bash]
$ cgitsync purge <spec.cgs> [--output-path DIR]
\end{lstlisting}

Removes generated clone state from the resolved \texttt{CGSHOME} without
cloning. It deletes child repository directories declared directly under the
project root, project-local \texttt{*.lgr} files, and the root
\texttt{.gitmodules} file. Nested directories that are not immediate root
children are left untouched by this command.

\subsection{\texttt{pull}}

\begin{lstlisting}[language=bash]
$ cgitsync pull [spec.cgs|snapshot.gts]
\end{lstlisting}

Resynchronises the project tree.  \texttt{.cgs} sources follow the
\texttt{restart} flow (load, discover, checkout current branch, READY), while
\texttt{.gts} sources follow saved-state restore semantics equivalent to
\texttt{launch-release}.  On success, the tree ends in \texttt{READY}.

\subsection{\texttt{clone}}

\begin{lstlisting}[language=bash]
$ cgitsync clone <spec.cgs> [--target-dir DIR]
\end{lstlisting}

Clones the full repository tree declared by the \texttt{.cgs} file. The root
project directory defaults to \texttt{./<project-name>}. If that path already
exists and is not empty, \cgscli{} chooses the next available suffixed
directory such as \texttt{./ComplexGitSync-1}. When \texttt{--target-dir} is
provided, that path is used exactly.

For each repository, \cgscli{} tries the repo target branch first
(typically \texttt{default\_branch = "autoTest"}) and falls back to
\texttt{fallback\_branch} (typically \texttt{"main"}) when the preferred
branch does not exist on the remote. Nested \texttt{.cgs} files are discovered
after parent repositories are cloned, so explicit and automatic nested trees
are materialised during the same command.

On success, \texttt{clone} exits~0 and prints a READY summary including the
selected root path. It also writes a runtime snapshot to
\texttt{<project-root>/.cgitsync/state/<spec-name>.gts} and registers that
snapshot as the latest runtime state for the source \texttt{.cgs} file.

\subsection{\texttt{checkout}}

\begin{lstlisting}[language=bash]
$ cgitsync checkout <branch> [--gts <snapshot.gts>] [--ref-kind branch|tag]
\end{lstlisting}

Checks out a branch or tag across the whole tree, parent-first.
Loads the READY registry from \texttt{--gts}, then propagates the ref,
creates missing local branches, and runs \texttt{git checkout} in every repo.
On success the tree remains \texttt{READY} and a new \texttt{.gts} snapshot is written.

\subsection{\texttt{add}}

\begin{lstlisting}[language=bash]
$ cgitsync add [--gts <snapshot.gts>] [--dry-run]
\end{lstlisting}

Stages all changes across the tree (\texttt{git add --all}), leaf-first.
Requires \texttt{READY}; tree stays \texttt{READY}. With
\texttt{--dry-run}, \cgscli{} prints the leaf-first execution plan and does not
mutate repositories.

\subsection{\texttt{commit}}

\begin{lstlisting}[language=bash]
$ cgitsync commit <message> [--gts <snapshot.gts>] [--no-stage] [--dry-run]
\end{lstlisting}

Commits changes across the tree, leaf-first.  Loads the READY registry from
\texttt{--gts}.  Repositories with no staged changes after the optional
\texttt{git add --all} step (suppressed by \texttt{--no-stage}) are silently
skipped. Requires a \texttt{READY} tree; the tree remains \texttt{READY} on success.
Commit preflight now emits warnings for actionable workspace issues such as
dirty worktrees, ahead branches, stale recorded \texttt{commit\_sha} values,
or child repositories that are not linked as submodules.
With \texttt{--dry-run}, \cgscli{} prints the planned stage/commit sequence
without performing git writes.

The commit message conventionally ends with \texttt{CGS\#VERSION}.

\subsection{\texttt{push}}

\begin{lstlisting}[language=bash]
$ cgitsync push [--gts <snapshot.gts>] [--dry-run]
\end{lstlisting}

Pushes all repositories to their configured remotes, leaf-first.  Loads the
READY registry from \texttt{--gts}.  Each repo is pushed to
\texttt{entry.remote\_name} (defaulting to \texttt{"origin"}) on its currently
resolved branch.  Refreshes the stored hash in \texttt{GitTree}.  Requires a
\texttt{READY} tree; the tree remains \texttt{READY} on success. Push preflight
blocks detached HEADs, missing remotes, unresolved merges, behind/diverged
branches, and missing submodule links; it warns when the local branch is ahead
or when the loaded snapshot records an older \texttt{commit\_sha}.
With \texttt{--dry-run}, \cgscli{} prints the push plan without mutating any
repository.

\subsection{\texttt{tag}}

\begin{lstlisting}[language=bash]
$ cgitsync tag <name> [--gts <snapshot.gts>] [--dry-run]
\end{lstlisting}

Creates and pushes the same tag across the full tree, leaf-first.  Loads the
READY registry from \texttt{--gts}.  Refreshes the stored tag in
\texttt{GitTree}.  Requires a \texttt{READY} tree and keeps the tree
\texttt{READY} on success. Tag preflight requires clean worktrees, aligned
branches, existing remotes, no pre-existing tag of the same name, no detached
HEADs or unresolved merges, and parent/child repositories linked as Git
submodules. Stale recorded \texttt{commit\_sha} values are reported as
warnings. Tag creation is non-forcing by default (no implicit \texttt{-f}).
With \texttt{--dry-run}, \cgscli{} prints the tag/push plan and exits without
creating tags.

\subsection{\texttt{freeze}}

\begin{lstlisting}[language=bash]
$ cgitsync freeze <name> [--gts <snapshot.gts>] [--dry-run]
\end{lstlisting}

Performs a release freeze across all repos, leaf-first: stage/commit,
create+push the shared release tag, refresh SHA/state, and write a
\texttt{.gts} snapshot. Loads the READY registry from \texttt{--gts}.
Requires \texttt{READY} and leaves the tree \texttt{READY}. Freeze preflight
checks remote availability, branch alignment, detached HEADs, unresolved
merges, tag absence, and submodule links for all parent/child repository
boundaries. Dirty worktrees and stale recorded \texttt{commit\_sha} values are
reported as warnings because freeze will stage, commit, and snapshot them.
Freeze snapshots also include a deterministic \texttt{freeze\_manifest} section
that records freeze invariants (immutable snapshot, validated workspace,
synchronized tag, ledger checkpoint) and declares \texttt{launch\_state} as the
canonical restore operation.
With \texttt{--dry-run}, \cgscli{} prints the planned add/commit/tag/push graph
for the loaded workspace without mutating it.

\texttt{freeze-release} and \texttt{freeze-state} remain available as
compatibility aliases for release and internal-state freeze operations
respectively.

The integration suite validates local no-network lifecycle coverage through
the CLI path, including initialisation and the READY git command cycle
\texttt{add -> commit -> push -> tag -> freeze}.

\subsection{\texttt{print}}

\begin{lstlisting}[language=bash]
$ cgitsync print <spec.cgs|snapshot.gts>
\end{lstlisting}

Prints a lifecycle summary as pretty JSON. For \texttt{.gts} input, it reports
the recorded snapshot state. For \texttt{.cgs} input, it reports the current
declared/runtime summary.

\subsection{\texttt{view-tree}}

\begin{lstlisting}[language=bash]
$ cgitsync view-tree [spec.cgs|snapshot.gts] [--depth N] [--collapse REPO]
\end{lstlisting}

Renders a deterministic terminal tree that focuses on repository topology and
operational state:
\texttt{name [branch] local\_state sync\_state}. Use \texttt{--depth} to
limit visible levels and repeat \texttt{--collapse} to hide selected
subtrees.

\subsection{\texttt{view-operation}}

\begin{lstlisting}[language=bash]
$ cgitsync view-operation [spec.cgs|snapshot.gts]
$ cgitsync view_operation [spec.cgs|snapshot.gts]
\end{lstlisting}

Renders the runtime operation table with required columns:
\texttt{REPOSITORY}, \texttt{BRANCH}, \texttt{LOCAL\_STATE}, and
\texttt{SYNC\_STATE}.

\subsection{\texttt{describe}}

\begin{lstlisting}[language=bash]
$ cgitsync describe <snapshot.gts>
\end{lstlisting}

Reads a \texttt{.gts} snapshot and prints its contents as pretty-printed
JSON. This is a compatibility alias of \texttt{print}; both commands share the
same implementation.

\subsection{\texttt{restart}}

\begin{lstlisting}[language=bash]
$ cgitsync restart <spec.cgs>
\end{lstlisting}

Resynchronizes an already-cloned project tree. Compatibility alias for
\texttt{pull} with a \texttt{.cgs} source.

\subsection{\texttt{launch-release}}

\begin{lstlisting}[language=bash]
$ cgitsync launch-release <snapshot.gts>
\end{lstlisting}

Loads a release \texttt{.gts}, reconstructs the tree, performs required
clone/checkout actions, and must end in \texttt{READY} or fail explicitly.

\subsection{\texttt{launch-state}}

\begin{lstlisting}[language=bash]
$ cgitsync launch-state <snapshot.gts>
\end{lstlisting}

Loads an internal-state \texttt{.gts}, reconstructs the tree, performs required
clone/checkout actions, and must end in \texttt{READY} or fail explicitly.

\subsection{\texttt{validate-topology}}

\begin{lstlisting}[language=bash]
$ cgitsync validate-topology --gts <snapshot.gts>
\end{lstlisting}

Inspects the workspace branch topology and prints a human-readable report:
each repository's active branch, the reference branch (root's branch), and any
detected conflicts.  Returns exit code~0 when the topology is coherent and~1
when blocking conflicts are present.

\textbf{Branch Topology Propagation Rules (T35):}

\begin{enumerate}
  \item \textbf{Reference branch}: The root repository's current branch is
        the canonical reference for all repos.
  \item \textbf{Leaf-to-root inheritance}: Branch targeting flows
        root-first via \texttt{propagate\_global\_branch}.
        \texttt{validate-topology} verifies the on-disk state without
        issuing any git writes.
  \item \textbf{Allowed divergence}: Repositories in a frozen
        (\texttt{TAG}) state produce a \texttt{tag\_divergence} diagnostic
        that does \emph{not} make the topology incoherent.
  \item \textbf{Incoherent states} (blocking): \texttt{misaligned\_branch}
        (different branch from root) and \texttt{detached\_head} (detached
        HEAD without a tag reference).
\end{enumerate}

\subsection{Planned Commands (not yet implemented)}

The following commands are defined in the CLI command table but are not yet
implemented. Each returns exit code~2 with the message
\texttt{"not implemented yet"}.

\begin{itemize}
  \item \texttt{write-gts} — write a \texttt{.gts} snapshot explicitly
    from a loaded \texttt{.cgs} or runtime state.
\end{itemize}

\subsection{\texttt{status}}

\begin{lstlisting}[language=bash]
$ cgitsync status [--gts FILE]
\end{lstlisting}

Summarises tree readiness and per-repository sync state. The table reports the
local branch, upstream branch, local worktree state, ahead/behind tracking
counts, current HEAD, and the commit recorded in the loaded \texttt{.gts}.

% -----------------------------------------------------------------------
\section{Document Formats}

\texttt{.cgs} means \textbf{ComplexGitSync} specification, \texttt{.gts} means
\textbf{GitTreeState} snapshot, \texttt{.goc} means
\textbf{GitOrchestrationCommand} plan, and \texttt{.lgr} means
\textbf{Local Git Register}.

\subsection{Local Authoring Spec (\texttt{.cgs})}

A TOML file that describes the repository tree to manage. It is the only
hand-authored input; it is never modified at runtime.

\textbf{Top-level sections:}

\begin{itemize}
  \item \texttt{[document]} — mandatory; contains \texttt{format\_version}.
  \item \texttt{[project]} — mandatory; project-level metadata (see
        Table~\ref{tab:project-fields}).
  \item \texttt{[[repos]]} — one entry per managed repository
        (see Table~\ref{tab:repo-fields}).
\end{itemize}

\begin{table}[h]
\centering
\begin{tabular}{lll}
\hline
\textbf{Field} & \textbf{Type} & \textbf{Description} \\
\hline
\texttt{name}           & string  & Human-readable project label \\
\texttt{repo\_name}     & string  & Git repository slug \\
\texttt{default\_branch}& string  & Default branch (e.g.\ \texttt{main}) \\
\texttt{gitprovider}    & string  & Provider (\texttt{github}, \texttt{gitlab}, \texttt{custom}) \\
\texttt{gitprovider\_url}& string & Required for \texttt{custom} provider \\
\hline
\end{tabular}
\caption{Project-level fields in \texttt{.cgs}}
\label{tab:project-fields}
\end{table}

\begin{table}[h]
\centering
\begin{tabular}{lll}
\hline
\textbf{Field} & \textbf{Type} & \textbf{Description} \\
\hline
\texttt{gitprovider}         & string  & Provider override for this repo \\
\texttt{project\_owner\_name}& string  & GitHub owner / GitLab group \\
\texttt{project\_name}       & string  & Repository name on the provider \\
\texttt{relative\_path}      & string  & Path relative to project root \\
\texttt{gitprovider\_url}    & string  & Required for \texttt{custom} provider \\
\hline
\end{tabular}
\caption{Per-repository fields in \texttt{.cgs}}
\label{tab:repo-fields}
\end{table}

\subsection{Git Tree State Snapshot (\texttt{.gts})}

A TOML file generated automatically by bootstrapping and release operations.
It must never be hand-edited. It captures the exact commit SHAs and lifecycle
state of the tree at a point in time.

\textbf{Top-level sections:}
\begin{itemize}
  \item \texttt{[document]} — \texttt{format\_version}, \texttt{schema\_version},
        \texttt{generated\_at}, \texttt{command\_origin},
        \texttt{hash\_algorithm}, \texttt{snapshot\_hash}.
  \item \texttt{[project]} — project name and \texttt{root\_absolute\_path}.
  \item \texttt{[tree\_state]} — \texttt{lifecycle\_state}, \texttt{is\_ready},
        \texttt{registry\_complete}.
  \item \texttt{[[repo\_state]]} — one entry per repo; includes
        \texttt{commit\_sha}, \texttt{sync\_state}, and ref information.
\end{itemize}

\textbf{Why both version fields and hash metadata exist:}
\begin{itemize}
  \item \texttt{format\_version} identifies the broad document family and keeps
        long-range compatibility intent explicit.
  \item \texttt{schema\_version} identifies the exact field-level contract used
        by validation and canonical snapshot serialization.
  \item \texttt{snapshot\_hash} (with \texttt{hash\_algorithm = "sha256"})
        provides deterministic workspace identity: identical tree state produces
        identical digest even if volatile metadata (\texttt{generated\_at},
        \texttt{command\_origin}) differs.
\end{itemize}

\textbf{How this is used at runtime:}
\begin{itemize}
  \item During snapshot generation, ComplexGitSync writes
        \texttt{schema\_version}, \texttt{hash\_algorithm}, and canonical
        \texttt{snapshot\_hash}.
  \item During load/validation, \texttt{snapshot\_hash} is recomputed from the
        canonical payload and must match.
  \item The project-local \texttt{.lgr} register uses this canonical hash to
        deduplicate equivalent snapshots.
\end{itemize}

\subsection{Planning / GOC Documents (\texttt{.goc})}

Used for project identity and planning metadata. Supported in TOML, YAML, or
JSON. YAML support requires the optional \texttt{PyYAML} dependency.

\subsection{Local Git Register and Sync Ledger (\texttt{.lgr})}

A TOML file maintained per project at \texttt{<project-root>/<Project\_name>.lgr}.
It is updated automatically whenever a \texttt{.gts} snapshot is written.

\textbf{Sections:}
\begin{itemize}
  \item \texttt{[register]} — current snapshot pointer (\texttt{current\_snapshot\_id},
        \texttt{current\_snapshot\_hash}, \texttt{current\_snapshot\_path}).
  \item \texttt{[[snapshots]]} — list of stable \texttt{gts-XXXXXX} identifiers,
        one per unique workspace state, deduplicated by SHA-256 hash.
  \item \texttt{[[ledger]]} — append-only list of synchronisation events forming
        a DAG via \texttt{parent\_sync\_ids}.
\end{itemize}

\textbf{Ledger event schema:}

\begin{lstlisting}[language=bash]
[[ledger]]
sync_id         = "lgr-000001"
parent_sync_ids = []
operation       = "clone"
timestamp       = "2026-05-20T19:48:50.159Z"
actor           = "username"
workspace_hash  = "5f7c8a2d..."   % 64-character SHA-256 hex digest
gts_snapshot_id = "gts-000001"
affected_repos  = ["demo", "dep-a"]
\end{lstlisting}

The \texttt{workspace\_hash} field equals the canonical SHA-256 digest
(\texttt{document.snapshot\_hash}) of the associated \texttt{.gts} file,
creating a direct link between each ledger event and the immutable workspace
state it records.  Events are parents-before-children in DAG order.

% -----------------------------------------------------------------------
\section{Configuration and Environment}

\cgspkg{} uses project files for topology and does not require global topology
configuration. Snapshot discovery can use \texttt{CGSHOME}; when it is not set,
the \texttt{initialise(.cgs)} output-path default is \texttt{CGSPATH=../..}; later
commands can also walk upward from the current directory to find
\texttt{.cgitsync/state/}. The logging and runtime-state subsystems use the
user state directory by default:

\begin{itemize}
  \item logs: \texttt{~/.local/state/ComplexGitSync/logs/}
  \item latest runtime-state mapping:
        \texttt{~/.local/state/ComplexGitSync/runtime/}
\end{itemize}

If \texttt{XDG\_STATE\_HOME} is set, these directories are rooted there
instead.

% -----------------------------------------------------------------------
\section{Error Reference}

\begin{description}
  \item[\texttt{ValueError}] Raised when a required field is missing or a
    field value is invalid (e.g.\ unknown \texttt{gitprovider},
    missing \texttt{gitprovider\_url} for a \texttt{custom} provider).
  \item[\texttt{GitSyncError}] Raised when clone-time Git operations fail,
    for example because the remote branch cannot be resolved or the target
    directory is already occupied.
  \item[Exit code 1] General CLI failure (parse error, validation error).
  \item[Exit code 2] Command exists but is not yet implemented.
\end{description}

% -----------------------------------------------------------------------
\section{Troubleshooting}

\begin{description}
  \item[``\texttt{gitprovider\_url} is required for custom provider'']
    Add a \texttt{gitprovider\_url} field to the offending
    \texttt{[[repos]]} entry in your \texttt{.cgs} file.
  \item[\texttt{validate} exits non-zero despite a correct-looking file]
    Check that \texttt{[document]} contains \texttt{format\_version = "1.0"}
    and that every \texttt{[[repos]]} block includes
    \texttt{project\_owner\_name} and \texttt{project\_name}.
  \item[\texttt{Clone destination already exists and is not empty}]
    Re-run \texttt{clone} with \texttt{--target-dir} to force a clean location,
    or let \cgscli{} choose a suffixed default directory.
  \item[\texttt{No cloneable branch found for <repo>}]
    Confirm that the remote exposes either the repo's
    \texttt{default\_branch} or its \texttt{fallback\_branch}.
  \item[PyYAML import error when using \texttt{.goc} files]
    Add PyYAML in the Pixi environment: \texttt{pixi add pyyaml}.
\end{description}
